\documentclass[10pt,a4paper]{amsart}
\usepackage[margin=19mm]{geometry}
\usepackage{amsmath,amssymb,mathtools}
\usepackage[scr=rsfs,cal=euler]{mathalfa}
\usepackage{xfrac}
\usepackage[unicode,hidelinks,bookmarks=false]{hyperref}
\newcommand{\ie}{i.e.\ }
\newcommand{\cf}{cf.\ }
\newcommand{\resp}{respectively\ }
\newcommand{\Subsection}{\subsection}
\providecommand{\nobreakdash}{\nobreak\hbox{-}}
\setlength{\emergencystretch}{2em}
\multlinegap0pt
\usepackage{amscd}
\usepackage{indentfirst}
\usepackage{eufrak}
\usepackage{tikz}
\usepackage{ulem}
\newtheorem{theorem}{Theorem}
\newtheorem{lemma}[theorem]{Lemma}
\newtheorem{remark}[theorem]{Remark}
\newtheorem{corollary}[theorem]{Corollary}
\newtheorem{proposition}[theorem]{Proposition}
\newtheorem{example}[theorem]{Example}
\newtheorem{definition}[theorem]{Definition}
\def\proofname{\noindent {\sl Proof.}}
\begin{document}
\title[Algebraic Lattices Arising from Congruence Submodules in Sub]{Algebraic Lattices Arising from Congruence Submodules in Subfields of $p$-th Cyclotomic Fields}
\author{Trajano Pires da Nóbrega Neto, Antonio Aparecido de Andrade, Jéfferson Luiz Rocha Bastos, Robson Ricardo de Araujo, José Carmelo Interlando}
\date{}
\maketitle
\noindent\emph{Source and licence.} Algebraic Lattices Arising from Congruence Submodules in Subfields of $p$-th Cyclotomic Fields. Version arXiv:2608.17056v1 (CC BY 4.0). This material is distributed under the Creative Commons Attribution 4.0 International licence, http://creativecommons.org/licenses/by/4.0/, as recorded in the arXiv OAI licence metadata for this exact version and shown on the abstract page https://arxiv.org/abs/2608.17056v1. Full rights notice: licenses/arxiv-2608.17056-CC-BY-4.0.txt.\par\smallskip
\noindent\textit{Editorial scope.} Counted unit: the original \texttt{theorem} environment \texttt{prop-3} at source lines 327--330 together with its complete proof at lines 331--333; the delivered document keeps source lines 100--334 so that every referenced label and every citation resolves inside it. Native editable source is the paper's own concatenated TeX; the AMS wrapper is editorial formatting only. No publisher class, upstream script or unrelated artwork is redistributed.
\section{Subfields of the $p$-th cyclotomic field} \label{sec-2}

Let $\mathbb{L}=\mathbb{Q}(\zeta_p)$ be the $p$-th cyclotomic field, where $p$ is a prime number. It is well known that, for every divisor $s$ of $p-1$, there exists a unique subfield $\mathbb{K}\subseteq \mathbb{L}$
with $[\mathbb{K}:\mathbb{Q}]=s$, by the Galois correspondence. Moreover, the rings of algebraic integers of $\mathbb{L}$ and of its maximal real subfield $\mathbb{Q}(\zeta_p+\zeta_p^{-1})$ are
\[ \mathcal{O}_{\mathbb{L}}=\mathbb{Z}[\zeta_p] \quad\text{and}\quad \mathcal{O}_{\mathbb{Q}(\zeta_p+\zeta_p^{-1})} =\mathbb{Z}[\zeta_p+\zeta_p^{-1}], \]
respectively \cite{wash}. Consequently, \[ \{1,\zeta_p,\zeta_p^2,\ldots,\zeta_p^{p-2}\} \ \ \mbox{and} \ \ \{1,\zeta_p+\zeta_p^{-1},\ldots,\zeta_p^{\frac{p-3}{2}}+\zeta_p^{-\frac{p-3}{2}}\} \]
are integral bases of $\mathbb{L}$ and $\mathbb{Q}(\zeta_p+\zeta_p^{-1})$, respectively.

Let $\mathbb{K}$ be an arbitrary subfield of $\mathbb{L}$, and write \[ s=[\mathbb{K}:\mathbb{Q}] \ \ \mbox{and} \ \ r=[\mathbb{L}:\mathbb{K}]=\frac{p-1}{s}. \] Then both $r$ and $s$ divide $p-1$. The Galois group
$G=\operatorname{Gal}(\mathbb{L}/\mathbb{Q})$ is canonically isomorphic to the multiplicative group $\mathbb{Z}_p^{\star} =(\mathbb{Z}/p\mathbb{Z})^{\star}$,
which is cyclic of order $p-1$. Let $\bar{u}$ be a generator of $\mathbb{Z}_p^{\star}$, and let $\sigma_u$ denote the automorphism of $\mathbb{L}$ determined by $\sigma_u(\zeta_p)=\zeta_p^{\,u}$.
Then \[ G=\langle\sigma_u\rangle = \{\sigma_u^0,\sigma_u,\sigma_u^2,\ldots,\sigma_u^{p-2}\}. \] Since every subgroup of a cyclic group is cyclic, both extensions $\mathbb{Q}\subseteq\mathbb{K}$ and $\mathbb{K}\subseteq\mathbb{L}$ are cyclic Galois extensions.

\begin{proposition} If $H=\operatorname{Gal}(\mathbb{L}/\mathbb{K}) = \langle \sigma_{u}^s\rangle$, then $\mathbb{K}=\mathbb{Q}(t)$, where
\[ t = \operatorname{Tr}_{\mathbb{L}/\mathbb{K}}(\zeta_p) = \sum_{j=0}^{r-1}\zeta_p^{\,u^{js}}. \]
\end{proposition}
\begin{proof} Since $H=\operatorname{Gal}(\mathbb{L}/\mathbb{K})$, the trace \[ t = \operatorname{Tr}_{\mathbb{L}/\mathbb{K}}(\zeta_p) = \sum_{\tau\in H}\tau(\zeta_p) \]
is fixed by every element of $H$. Hence $t\in\mathbb{L}^H=\mathbb{K}$. Let $ G=\operatorname{Gal}(\mathbb{L}/\mathbb{Q}) = \langle\sigma_u\rangle$, where $\sigma_u(\zeta_p)=\zeta_p^{\,u}$. Since
$H=\langle\sigma_{u}^s\rangle$,  its elements are \[ H = \{\sigma_u^0,\sigma_{u}^s,\sigma_{u}^{2s},\ldots,\sigma_{u}^{(r-1)s}\}. \] Therefore, $t = \zeta_p+\zeta_p^{u^s}
+\zeta_p^{u^{2s}} +\cdots+\zeta_p^{u^{(r-1)s}}$. For each integer $i$, with $0\le i\le s-1$, \[ \sigma_u^{\,i}(t) = \sum_{j=0}^{r-1} \zeta_p^{\,u^{\,i+js}}. \]
The exponents $u^{i},u^{i+s},\ldots,u^{i+(r-1)s}$ are precisely the elements of the coset $u^{i}H$ of $H$ in $G$. Since the cosets $H,uH,\ldots,u^{s-1}H$ are distinct, the conjugates
$t,\sigma_u(t),\ldots,\sigma_u^{\,s-1}(t)$ are pairwise distinct. Hence the orbit of $t$ under $G$ has cardinality $s$, so that
$[\mathbb{Q}(t):\mathbb{Q}] =s$. Since $t\in\mathbb{K}$, it follows that $\mathbb{Q}(t)\subseteq\mathbb{K}$. Finally, $
[\mathbb{Q}(t):\mathbb{Q}] = [\mathbb{K}:\mathbb{Q}] = s$, which implies $\mathbb{Q}(t)=\mathbb{K}$. Thus 
\[ t = \operatorname{Tr}_{\mathbb{L}/\mathbb{K}}(\zeta_p) = \sum_{j=0}^{r-1}\zeta_p^{\,u^{js}} \] is a primitive generator of the field $\mathbb{K}$.
\end{proof}

\begin{proposition} \label{lem-1} If $\zeta_p^{\,u^{i_1s+j_1}} = \zeta_p^{\,u^{i_2s+j_2}}$, where $1\le i_1,i_2\le r$ and $1\le j_1,j_2\le s$, then $i_1=i_2$ and $j_1=j_2$. \end{proposition}
\begin{proof} For hypothesis, $\zeta_p^{\,u^{i_1s+j_1}} = \zeta_p^{\,u^{i_2s+j_2}}\iff u^{i_1s+j_1} \equiv u^{i_2s+j_2} \pmod p \iff i_1s+j_1 \equiv i_2s+j_2 \pmod{p-1}$. Thus, $(i_1-i_2)s=(j_2-j_1)+rst$, for some $t\in\mathbb{Z}$. Therefore, $j_2\equiv j_1 \pmod s$, that is, $j_1=j_2$. So, $(i_1-i_2)s=rst$ and $i_1\equiv i_2 \pmod r$. Hence, $i_1=i_2$. \end{proof}

\begin{proposition}\label{prop-1} The set $\left\{ \zeta_p^{\,u^{is+j}}:\; 0\le i\le r-1,\; 1\le j\le s \right\}$  is linearly independent over $\mathbb{Q}$. \end{proposition}
\begin{proof} Suppose that \[ \sum_{j=1}^{s}\sum_{i=0}^{r-1} a_{ji}\,\zeta_p^{\,u^{is+j}} = 0, \] where $a_{ji}\in\mathbb{Q}$.
By Proposition \ref{lem-1}, the exponents $u^{is+j}$, with $0\le i\le r-1$ and $1\le j\le s$, are pairwise distinct modulo $p$. Thus, there are $rs=p-1$
such exponents. Hence this relation can be rewritten as \[ \sum_{k=1}^{p-1} b_k\zeta_p^{\,k} = 0, \] where each $b_k\in\mathbb{Q}$.
Since $\{\zeta_p,\zeta_p^2,\ldots,\zeta_p^{p-1}\}$ is a $\mathbb{Q}$-basis of $\mathbb{Q}(\zeta_p)$, it follows that $b_k=0$, for $1\le k\le p-1$. 
Therefore, \[ a_{ji}=0, \ \ \mbox{for} \ \ 0\le i\le r-1 \ \ \mbox{and} \ \ \;1\le j\le s), \] which proves that the given set is linearly independent over $\mathbb{Q}$. \end{proof}

Denote $
t=\operatorname{Tr}_{\mathbb{L}/\mathbb{K}}(\zeta_p)$. Since $t$ is fixed by every element of $H=\operatorname{Gal}(\mathbb{L}/\mathbb{K})$, it follows that $ t\in\mathbb{K}$, and therefore,
\[ \mathbb{Q}(t)\subseteq\mathbb{K}. \] Moreover, $t$ is an algebraic integer because it is the trace of the algebraic
integer $\zeta_p$. Hence $\operatorname{Tr}_{\mathbb{Q}(t)/\mathbb{Q}}(t)\in\mathbb{Z}$. Using the transitivity of the trace, it follows that
\[  -1 = \operatorname{Tr}_{\mathbb{L}/\mathbb{Q}}(\zeta_p) = \operatorname{Tr}_{\mathbb{K}/\mathbb{Q}} \!\left( \operatorname{Tr}_{\mathbb{L}/\mathbb{K}}(\zeta_p)\right) = \operatorname{Tr}_{\mathbb{K}/\mathbb{Q}}(t) = [\mathbb{K}:\mathbb{Q}(t)] \operatorname{Tr}_{\mathbb{Q}(t)/\mathbb{Q}}(t). \]
Since the right-hand side is the product of the positive integer $[\mathbb{K}:\mathbb{Q}(t)]$ and an integer, it follows that $[\mathbb{K}:\mathbb{Q}(t)]=1$.
Therefore, $\mathbb{K}=\mathbb{Q}(t)$.

\begin{theorem}\label{teo-1} The set $S = \left\{t,\,\sigma_u(t),\,\sigma_u^{2}(t),\,\ldots,\,\sigma_u^{\,s-1}(t)\right\}$ is a $\mathbb{Q}$-basis of $\mathbb{K}$. \end{theorem}
\begin{proof} Since $\operatorname{Gal}(\mathbb{L}/\mathbb{Q}) = \langle\sigma_u\rangle$ is cyclic of order $p-1$, the subgroup $H=\operatorname{Gal}(\mathbb{L}/\mathbb{K})$
has order $r=(p-1)/s$ and is generated by $\sigma_{u^s}$. Thus $H = \left\{ \sigma_u^0,\sigma_{u}^s,\sigma_{u}^{2s},\ldots,\sigma_{u}^{(r-1)s}\right\}$. Consequently, \[ t = \operatorname{Tr}_{\mathbb{L}/\mathbb{K}}(\zeta_p) = \sum_{\tau\in H}\tau(\zeta_p) =
\zeta_p+\zeta_p^{u^s}+\zeta_p^{u^{2s}}+\cdots+\zeta_p^{u^{(r-1)s}}. \] Since $\mathbb{K}/\mathbb{Q}$ is Galois, every automorphism $\sigma_u^i$ restricts to an automorphism of $\mathbb{K}$. Therefore,
$\sigma_u^i(t)\in\mathbb{K}$, for $0\le i\le s-1$. Moreover, \[ \sigma_u^i(t) = \sum_{j=0}^{r-1} \zeta_p^{\,u^{\,js+i}}, \ \ \mbox{for} \ \ 0\le i\le s-1. \]
Suppose that \[ \sum_{i=0}^{s-1} a_i\sigma_u^i(t)=0, \ \ \mbox{for} \ \ a_i\in\mathbb{Q}. \] Then \[ \sum_{i=0}^{s-1} a_i \left(
\sum_{j=0}^{r-1} \zeta_p^{\,u^{\,js+i}} \right) = 0. \] By Proposition~\ref{prop-1}, the set \[ \left\{ \zeta_p^{\,u^{\,js+i}}: 0\le i\le s-1,\; 0\le j\le r-1\right\} \]
is linearly independent over $\mathbb{Q}$. Hence $a_i=0$, for $0\le i\le s-1$. Therefore, \[ S = \left\{ t,\sigma_u(t),\ldots,\sigma_u^{\,s-1}(t) \right\} \] is linearly independent over $\mathbb{Q}$.
Finally, $|S| = s = [\mathbb{K}:\mathbb{Q}]$, so $S$ is a $\mathbb{Q}$-basis of $\mathbb{K}$. \end{proof}

\begin{corollary} \label{cor-1} The set $S=\left\{t,\,\sigma_u(t),\,\sigma_u^{2}(t),\,\ldots,\,\sigma_u^{\,s-1}(t)\right\}$ is a $\mathbb{Z}$-basis of $\mathcal{O}_{\mathbb{K}}$. \end{corollary}
\begin{proof} By Theorem~\ref{teo-1}, the set \[ S = \left\{ t,\,\sigma_u(t),\,\ldots,\,\sigma_u^{\,s-1}(t)\right\} \] is a $\mathbb{Q}$-basis of $\mathbb{K}$.
First, we prove that $\langle S\rangle_{\mathbb{Z}}\subseteq\mathcal{O}_{\mathbb{K}}$. Let \[ \alpha = \sum_{i=0}^{s-1}a_i\sigma_u^i(t), \ \ \mbox{with} \ \  a_i\in\mathbb{Z}. \]
Since \[ t = \sum_{j=0}^{r-1} \zeta_p^{\,u^{js}}, \] it follows that \[ \sigma_u^i(t) = \sum_{j=0}^{r-1} \zeta_p^{\,u^{js+i}}, \ \ \mbox{for} \ \ 0\le i\le s-1. \]
Each $\sigma_u^i(t)$ is a sum of roots of unity and therefore belongs to
$\mathbb{Z}[\zeta_p]=\mathcal{O}_{\mathbb{L}}$. Since $\sigma_u^i(t)\in\mathbb{K}$ and $\mathcal{O}_{\mathbb{K}} = \mathcal{O}_{\mathbb{L}}\cap\mathbb{K}$, it follows that
\[ \sigma_u^i(t)\in\mathcal{O}_{\mathbb{K}}. \] Hence $\alpha\in\mathcal{O}_{\mathbb{K}}$, and consequently $\langle S\rangle_{\mathbb{Z}} \subseteq \mathcal{O}_{\mathbb{K}}$.
Conversely, let $\alpha\in\mathcal{O}_{\mathbb{K}}$. Since $S$ is a $\mathbb{Q}$-basis of $\mathbb{K}$, there exist unique
coefficients $a_0,\ldots,a_{s-1}\in\mathbb{Q}$ such that \[ \alpha = \sum_{i=0}^{s-1} a_i\sigma_u^i(t). \] Using the above expression for the conjugates of $t$, it follows that 
\[ \alpha = \sum_{i=0}^{s-1} a_i \left( \sum_{j=0}^{r-1} \zeta_p^{\,u^{js+i}} \right). \] By Proposition \ref{lem-1}, the exponents $u^{js+i}$, for $0\le i\le s-1$ and $0\le j\le r-1$,
are pairwise distinct modulo $p$. Since there are $rs=p-1$ such exponents, they are precisely the nonzero residue classes modulo $p$. Hence
\[ \alpha = \sum_{k=1}^{p-1} b_k\zeta_p^k, \] where each coefficient $b_k$ is one of the numbers $a_0,\ldots,a_{s-1}$. Now, \[
\alpha\in\mathcal{O}_{\mathbb{K}} \subseteq \mathcal{O}_{\mathbb{L}} = \mathbb{Z}[\zeta_p]. \] Since $\{\zeta_p,\zeta_p^2,\ldots,\zeta_p^{p-1}\}$ is a $\mathbb{Z}$-basis of $\mathbb{Z}[\zeta_p]$, it follows that $b_k\in\mathbb{Z}$, for $1\le k\le p-1$. 
Therefore, $a_i\in\mathbb{Z}$, for $0\le i\le s-1$. Hence $\alpha\in\langle S\rangle_{\mathbb{Z}}$,  which proves that $\mathcal{O}_{\mathbb{K}} \subseteq\langle S\rangle_{\mathbb{Z}}$.
So $\mathcal{O}_{\mathbb{K}} = \langle S\rangle_{\mathbb{Z}}$, and therefore, \[ S= \left\{t,\,\sigma_u(t),\,\sigma_u^{2}(t),\, \ldots,\,\sigma_u^{\,s-1}(t) \right\} \] is a $\mathbb{Z}$-basis of $\mathcal{O}_{\mathbb{K}}$.
\end{proof}

\begin{remark} \label{rem-0} A cyclic number field $\mathbb{K}$ is said to admit an integral normal basis if $\{t,\theta(t),\theta^2(t),\ldots,\theta^{s-1}(t)\}$ is a $\mathbb{Z}$-basis for $\mathcal{O}_{\mathbb{K}}$, where $t\in \mathcal{O}_\mathbb{K}$ and $\theta$ is a generator of $\textrm{Gal}(\mathbb{K}/\mathbb{Q})$. Hence, from Corollary \ref{cor-1}, every subfield of $\mathbb{Q}(\zeta_p)$ admits an integral normal basis given by $\{t,\theta(t),\theta^2(t),\ldots,\theta^{s-1}(t)\}$, where $\theta=\sigma_{u}|_\mathbb{K}$. \end{remark}

\section{The trace form} \label{sec-3}

%some applications of quadratic forms to cyclotomic subfields $\mathbb{K}$ of $\mathbb{Q}\zeta_p)$, where $p$ is a prime number. Our objective is to find 
Let $\mathbb{L}=\mathbb{Q}(\zeta_p)$ be a $p$-th cyclotomic field, where $p$ is a prime number. Let $\mathbb{K}$ be a number field such that $\mathbb{K}\subseteq\mathbb{L}$, with $[\mathbb{K}:\mathbb{Q}]=s$ and $[\mathbb{L}:\mathbb{K}]=r$. Let $\theta$ be the generator of the Galois group $Gal(\mathbb{K}/\mathbb{Q})$. Let $\mathcal{O}_{\mathbb{L}}=\mathbb{Z}[\zeta_p]$ and $\mathcal{O}_{\mathbb{K}}$ be the ring of algebraic integers of $\mathbb{L}$ and $\mathbb{K}$, respectively.
In this section, we present an expression of the integral trace form, i.e., the trace of an element of the form $x\overline{x}$, where $x$ is an algebraic integer of $\mathcal{O}_{\mathbb{K}}$. The minimum of this trace form is an important parameter in the study of the packing density of the lattice realized by $\mathcal{O}_\mathbb{K}$ via the canonical embedding.

%Let $p$ be a odd prime and $\mathbb{L}=\mathbb{Q}(\zeta_p)$. %$\mathcal{O}_{\mathbb{L}}=\mathbb{Z}[\zeta_p]$ be the ring of algebraic integers of $\mathbb{L}$ and $x=\displaystyle \sum_{i=0}^{p-2}a_i\zeta_p^i\in\mathbb{Z}[\zeta_p]$. Thus, 
It is well-known that $$Tr_{\mathbb{L}}(\zeta_p^k)=\begin{cases} -1 & \textrm{if } k\not\equiv 0 \pmod p,\\ p-1 & \textrm{if } k\equiv 0\pmod p. \end{cases}$$ If $x=\displaystyle \sum_{i=0}^{p-2}a_i\zeta_p^i\in\mathcal{O}_\mathbb{L}=\mathbb{Z}[\zeta_p]$, with $a_0,a_1,\ldots,a_{p-2}\in\mathbb{Z}$, then
%\[ \begin{array}{lll} Tr_{\mathbb{L}}(x) &=& \displaystyle \sum_{i=0}^{p-2}Tr_{\mathbb{L}}\left(a_i\zeta_p^i\right) = \displaystyle Tr_{\mathbb{L}}(a_0)+\sum_{i=1}^{p-2}Tr_{\mathbb{L}}\left(a_i\zeta_p^i\right)\\ &=& \displaystyle (p-1)a_0-\sum_{i=1}^{p-2}a_i=pa_0-\sum_{i=0}^{p-2}a_i. \end{array} \] 
%Since $\overline{\zeta_p}=\zeta_p^{-1}$, it follows that $\overline{x}=\displaystyle\sum_{i=0}^{p-2}a_i\zeta_p^{-i}$. Thus,
\[  \begin{array}{lll} x\overline{x} &=& (a_0+a_1\zeta_p+\cdots+a_{p-2}\zeta_p^{p-2})(a_0+a_1\zeta_p^{-1}+\cdots+a_{p-2}\zeta_p^{-(p-2)})\vspace{.2cm} \\ &=&(a_0^2+\cdots +a_{p-2}^2)+(a_0a_1+\cdots +a_{p-3}a_{p-2})(\zeta_p+\zeta_p^{-1})\vspace{.2cm}\\ && \displaystyle +\cdots + (a_0a_{p-3}+a_1a_{p-2})(\zeta_p^{p-3}+\zeta_p^{-(p-3)})+a_0a_{p-2}(\zeta_p^{p-2}+\zeta_p^{-(p-2)}).\end{array} \]
Denote $x_i=\zeta_p^{i}+\zeta_p^{-i}$ and $A_i=a_0a_i+a_1a_{i+1}+\cdots +a_{p-2-i}a_{p-2}$, for $i=1,\ldots,p-2$. With this notation, \begin{equation} \label{expxxbarra} x\overline{x}=A_0+A_1x_1+\cdots +A_{p-2}x_{p-2}. \end{equation}
Consequently, the trace form associated to $\mathbb{L}$ is given by the following:

\begin{theorem} \label{teo-2} If $x=\displaystyle\sum_{i=0}^{p-2}a_i\zeta_p^i\in\mathbb{Z}[\zeta_p]$, then \[ Tr_{\mathbb{L}}(x\overline{x})=\displaystyle p\sum_{i=0}^{p-2}a_i^2 - \left(\sum_{i=0}^{p-2}a_i\right)^2. \]\end{theorem} \begin{proof} %Consider $x=\displaystyle\sum_{i=0}^{p-2}a_i\zeta_p^i\in\mathbb{Z}[\zeta_p]$, with $a_i \in \mathbb{Z}$, for $i=1,2,\ldots,p-2$.
From Equation (\ref{expxxbarra}), $x\overline{x}=A_0+A_1x_1+\cdots + A_{p-2}x_{p-2}$, where $A_i=a_0a_i+a_1a_{i+1}+\cdots +a_{p-2-i}a_{p-2}.$ Since
\[ Tr_{\mathbb{L}}(x_i)=Tr_{\mathbb{L}}(\zeta_p^i+\zeta_p^{-i}) = Tr_{\mathbb{L}}(\zeta_p^i)+Tr_{\mathbb{L}}(\zeta_p^{-i})=-1-1=-2\]
and the coefficient $x_0$ of $A_0$ is equal to $1$, we have that \[ Tr_{\mathbb{L}}(1)=[\mathbb{L}:\mathbb{Q}]\cdot 1 = p-1. \] Thus,
\[ \begin{array}{lll} Tr_{\mathbb{L}}(x\overline{x}) &=& (p-1)A_0-2(A_1+A_2+\cdots+A_{p-2}) =(p-1)A_0\vspace{.2cm}\\ && -2(a_0a_1+\cdots+a_{p-3}a_{p-2}+a_0a_2+\cdots +a_{p-4}a_{p-2}+\cdots\vspace{.2cm}\\
&& +a_0a_{p-3}+a_1a_{p-2}+a_0a_{p-2}).\end{array} \] Hence, \[ Tr_{\mathbb{L}}(x\overline{x})=p\displaystyle\sum_{i=0}^{p-2}a_i^2-\left(\displaystyle\sum_{i=0}^{p-2}a_i^2+2\displaystyle\sum_{0\leq
i< j\leq p-2}a_ia_j\right). \] Therefore, \begin{equation} \label{eqtra1} Tr_{\mathbb{L}}(x\overline{x})=p\displaystyle\sum_{i=0}^{p-2}a_i^2-\left(\displaystyle\sum_{i=0}^{p-2}a_i\right)^2,
\end{equation} which proves the result. \end{proof}

From Equation ($\ref{eqtra1}$), it follows that \[ \begin{array}{lll} \displaystyle p\sum_{i=0}^{p-2}a_i^2-\left(\sum_{i=0}^{p-2}a_i\right)^2 & = &
 %= p(a_0^2+\cdots+a_{p-2}^2)-((a_0^2+\cdots+a_{p-2}^2)+a_0a_1+\cdots+a_{p-3}a_{p-2})\\ 
 \displaystyle (p-1)\sum_{i=0}^{p-2}a_i^2 - \displaystyle\sum_{o\leq i<j\leq p-2} a_ia_j \vspace{.2cm}\\ &=& \displaystyle\sum_{i=0}^{p-2}a_i^2
+(p-2)\sum_{i=0}^{p-2}a_i^2- \sum_{o\leq i<j\leq p-2}a_ia_j \vspace{.2cm} \\ &=&   \displaystyle \sum_{i=0}^{p-2}a_i^2 + \sum_{o\leq i<j\leq p-2}(a_i-a_j)^2. \end{array} \]
Therefore, \[ Tr_{\mathbb{L}}(x\overline{x})=\displaystyle\sum_{i=0}^{p-2}a_i^2+\displaystyle\sum_{0\leq i< j\leq p-2}(a_i-a_j)^2. \] Furthermore, from Equation (\ref{eqtra1}), \[ \min\{ Tr_{\mathbb{L}/\mathbb{Q}}(x\overline{x}); \  x \in \mathcal{O}_{\mathbb{L}}, \ x \neq 0\} = p-1, \] where the minimum is attained at $x=1$. %This minimization of the quadratic form will be used in the study of algebraic lattices, through the calculation of the center density of these lattices.

%Now, consider $\mathbb{K} \subseteq \mathbb{L}$ of degree $[\mathbb{K}:\mathbb{Q}]=s$ and denote $r=(p-1)/s=[\mathbb{L}:\mathbb{K}]$, $t=Tr_{\mathbb{L}:\mathbb{K}}(\zeta_p)$ and $\alpha \in \mathbb{Z}$ such that $\sigma_ {u}$ generates $G=\textrm{Gal}(\mathbb{L}/\mathbb{Q})$.

\begin{theorem}\label{prop500} If $x=\displaystyle \sum_{i=0}^{s-1} a_i\theta^i(t) \in \mathcal{O}_{\mathbb{K}}$, where $a_0,a_1,\ldots,a_{s-1}\in\mathbb{Z}$ and $t=Tr_{\mathbb{L}/\mathbb{K}}(\zeta_p)$, then \[ Tr_{\mathbb{K}}(x\overline{x}) = p\sum_{i=0}^{s-1}a_i^2-\frac{p-1}{s}\left(\sum_{i=0}^{s-1}a_i\right)^2 . \] \end{theorem} 
\begin{proof} Let $x=\displaystyle \sum_{i=0}^{s-1} a_i\theta^i(t) \in \mathcal{O}_{\mathbb{K}}$, with $a_0,a_1,\ldots,a_{s-1}\in\mathbb{Z}$. Since $t=Tr_{\mathbb{L}/\mathbb{K}}(\zeta_p)\in \mathcal{O}_{\mathbb{K}}$, it follows that $t={\zeta_p}^{u^s}+{\zeta_p}^{u^{2s}}+\cdots+{\zeta_p}^{u^{rs}}$. So, $\sigma_{u^i}(t)={\zeta_p}^{u^{s+i}}+\cdots+{\zeta_p}^{u^{rs+i}}$, for $i=0,1,\ldots,s-1$, and
$$x = \sum_{i=0}^{s-1} a_i\left(\sum_{j=0}^{r-1}\zeta_p^{u^{js+i}} \right),$$ that is, $x=a_0({\zeta_p}^{u^s}+{\zeta_p}^{u^{2s}}+\cdots+{\zeta_p}^{u^{rs}})+a_1({\zeta_p}^{u^{s+1}}+{\zeta_p}^{u^{2s+1}}+\cdots+{\zeta_p}^{u^{rs+1}})+\cdots+a_{s-1}({\zeta_p}^{u^{s+s-1}}+\cdots+{\zeta_p}^{u^{rs+s-1}})$. 
%\[ x=a_1({\zeta_p}^{u^{s+1}}+\cdots+{\zeta_p}^{u^{rs+1}})+a_2({\zeta_p}^{u^{s+2}}+\cdots+ {\zeta_p}^{u^{rs+2}})+\cdots+a_s({\zeta_p}^{u^{s+s}}+\cdots+{\zeta_p}^{u^{rs+s}}). \]
From Theorem \ref{teo-2}, it follows that \[ Tr_{\mathbb{L}}(x\overline{x}) = r\sum_{i=0}^{s-1}a{_{i}^{2}} + {r^2}\sum_{0\leq i < j \leq s-1}^{}{(a_i - a_j)}^{2}. \] Since  $Tr_{\mathbb{L}/\mathbb{Q}}(x\overline{x}) = [\mathbb{L}:\mathbb{K}] Tr_{\mathbb{K}}(x\overline{x})$, then
$$Tr_{\mathbb{K}}(x\overline{x})=\frac{1}{r}Tr_{\mathbb{L}/\mathbb{Q}}(x\overline{x}),$$ and so, \[ Tr_{\mathbb{K}}(x\overline{x})=\sum_{i=0}^{s-1}a_{i}^{2} + r\sum_{0 \leq i < j \leq s-1}^{}{(a_i - a_j)}^{2}. \] Therefore, \[ Tr_{\mathbb{K}}(x\overline{x}) = p\sum_{i=0}^{s-1}a_i^2-\frac{p-1}{s} \left(\sum_{i=0}^{s-1}a_i\right)^2, \] which proves the result. \end{proof} 

\section{Algebraic lattices} \label{sec-4}

If $\mathbb{K}$ is an algebraic number field of degree $n$, there are exactly $n$ distinct $\mathbb{Q}$-monomorphisms $\sigma_i:\mathbb{K}\rightarrow\mathbb{C}$, for $i=1,2,\ldots,n$. A $\mathbb{Q}$-monomorphism $\sigma_i$ is said to be real if $\sigma_i(\mathbb{K})\subseteq\mathbb{R}$. Otherwise, $\sigma_i$ is said to be imaginary. It is a fact that $n=r_1+2r_2$, where $r_1$ is the number of real $\mathbb{Q}$-monomorphisms of $\mathbb{K}$ and $r_2$ is the number of pairs of imaginary $\mathbb{Q}$-monomorphisms. The number field $\mathbb{K}$ is said to be totally real if $\sigma_i$ is real for all $i=1,2,\ldots,n$. In turn, the number field $\mathbb{K}$ is totally imaginary if $\sigma_i$ is imaginary for all $i=1,2,\ldots,n$. 

Let $\sigma_1,\sigma_2,\ldots,\sigma_n$ denote the $\mathbb{Q}$-monomorphisms of $\mathbb{K}$ ordered such that $\sigma_1,\ldots,\sigma_{r_1}$ are the real monomorphisms and $\sigma_{r_1+1},\ldots,\sigma_{r_1+r_2}$ and $\sigma_{r_1+r_2+j}=\overline{\sigma_{r_1+j}}$, for $j=1,2,\ldots,r_2$, are the imaginary monomorphisms. The trace of an element $x\in\mathbb{K}$ over $\mathbb{Q}$ can be computed by $Tr_{\mathbb{K}}(x)=\sum_{i=1}^n\sigma_i(x)$. The discriminant of $\mathbb{K}$ over $\mathbb{Q}$ is defined by $D(\mathbb{K})=\det(Tr_{\mathbb{K}}(\alpha_i\alpha_j))_{i,j=1}^n$, where $\{\alpha_1,\alpha_2,\ldots,\alpha_n\}$ is an integral basis of $\mathcal{O}_{\mathbb{K}}$ (for details, see, e.g., \cite{samuel}). 

The canonical embedding $\sigma:\mathbb{K}\rightarrow\mathbb{R}^{r_1}\times\mathbb{C}^{r_2}$ is defined by  \[ \sigma(x) = (\sigma_1(x),\ldots,\sigma_{r_1}(x),\sigma_{r_1+1}(x),\ldots,\sigma_{r_1+r_2}(x)), \] where $x\in\mathbb{K}$ \cite[Section 4.2]{samuel}. We shall frequently identify $\mathbb{R}^{r_1}\times\mathbb{C}^{r_2}$ with $\mathbb{R}^n$, and thus, we consider $\sigma:\mathbb{K}\rightarrow\mathbb{R}^{n}$ given, for each $x\in \mathbb{K}$, by %given by \[ \begin{array}{lll} \sigma_{\alpha}(x) &=& (\sqrt{\alpha_1}\sigma_1(x),\ldots,\sqrt{\alpha_{r_1}}\sigma_{r_1}(x),\Re(\sqrt{\alpha_{r_1+1}}\sigma_{r_1+1}(x)),\\ && \Im(\sqrt{\alpha_{r_1+1}}\sigma_{r_1+1}(x)),\ldots,\Re(\sqrt{\alpha_{r_1+r_2}}\sigma_{r_1+r_2}(x)),\Im(\sqrt{\alpha_{r_1+r_2}}\sigma_{r_1+r_2}(x)), \end{array} \] where $x\in\mathbb{K}$. 
\begin{multline*} \sigma(x) = (\sigma_1(x),\ldots,\sigma_{r_1}(x),\mathfrak{R}(\sigma_{r_1+1}(x)),\mathfrak{I}(\sigma_{r_1+1}(x)),\\  \ldots,\mathfrak{R}(\sigma_{r_1+r_2}(x)),\mathfrak{I}(\sigma_{r_1+r_2}(x))), \end{multline*}
where $\mathfrak{R}(\beta)$ and $\mathfrak{I}(\beta)$ denote the real and imaginary parts of the complex number $\beta$, respectively.

Suppose $\overline{\sigma_i(x)}=\sigma_i(\overline{x})$ for all $x\in\mathbb{K}$ and for all $i=0,1,\ldots,n$. If $\mathcal{M}$ is a free $\mathbb{Z}$-module of $\mathcal{O}_{\mathbb{K}}$ of rank $n$, then $\sigma(\mathcal{M})$ is an $n$-dimensional lattice whose minimum is given by $\min\{||\sigma(x)||^2:x\in\mathcal{M},\ x\neq 0\}$, where \[ ||\sigma(x)||^2 = \left\{ \begin{array}{ll} Tr_{\mathbb{K}}(x^2) & \mbox{if} \ \mathbb{K} \ \mbox{is totally real,}\\  \frac{1}{2}Tr_{\mathbb{K}}(x\overline{x}) & \mbox{if} \ \mathbb{K} \ \mbox{is totally complex}. \end{array} \right. \] %that is, the computation of the trace form $Tr_{\mathbb{m}}(x\overline{x})$ is an interesting problem in algebraic number theory that has an important connection with the theory of lattices.

The center density of the lattice $\sigma(\mathcal{M})$ is given by \begin{equation} \label{densidade} \delta(\sigma(\mathcal{M})) = \frac{\rho^{n/2}}{2^n\sqrt{|D(\mathbb{K})|}[\mathcal{O}_{\mathbb{K}}:\mathcal{M}]}, \end{equation} where $\rho$ denotes the minimum of trace form of $\mathbb{K}$ over $\mathcal{M}$ and $[\mathcal{O}_{\mathbb{K}}:\mathcal{M}]$ denotes the index of $\mathcal{M}$ in $\mathcal{O}_{\mathbb{K}}$ (see, e.g., \cite{samuel}). 
In particular, if $\mathbb{K}$ is a subfield of $\mathbb{L}=\mathbb{Q}(\zeta_p)$ of degree $s$, then $D(\mathbb{K})=\pm p^{s-1}$ (\cite[Corollary 4.2]{trajano}) and 
\begin{equation}\label{eq_delta} \delta(\sigma(\mathcal{M})) = \frac{\rho^{s/2}}{2^sp^{\frac{s-1}{2}}[\mathcal{O}_{\mathbb{K}}:\mathcal{M}]}. \end{equation}
Furthermore, the map \[ \begin{array}{llll} \varphi:&\mathbb{Z}^s & \rightarrow & \mathcal{O}_{\mathbb{K}}\\ & {\bf x}=(a_0,a_1,\ldots,a_{s-1}) & \mapsto & \displaystyle \varphi({\bf x})=\sum_{i=0}^{s-1}a_i\theta^i(t), \end{array} \] is an isomorphism of $\mathbb{Z}$-modules, where $\theta$ is a generator of $Gal(\mathbb{K}/\mathbb{Q})$. The map $\phi=Tr_{\mathbb{K}}\circ\varphi:\mathbb{Z}^s\rightarrow \mathbb{N}\cup\{0\}$ given by \[ \phi({\bf x})=p\sum_{i=0}^{s-1}a_i^2-\frac{p-1}{s}\left(\sum_{i=0}^{s-1}a_i\right)^2, \] for each ${\bf x}=(a_0,a_1,\ldots,a_{s-1})\in\mathbb{Z}^s$, is a positive-definite quadratic form.

%Let \[ \mathcal{M}_m = \left\{x=\sum_{i=0}^{s-1}a_i\theta^i(t):a_0+a_1+\cdots+a_{s-1}\equiv 0(mod\ m)\right\}, \] where $Gal(\mathbb{Q}(t)/\mathbb{Q})=<\theta>$. In this case, $[\mathcal{O}_{\mathbb{K}}:\mathcal{M}_m]=m$.

\section{Constructions of algebraic lattices} \label{sec-5}

In this section, we provide a construction of algebraic lattices coming from certain modules in the ring of integers of $\mathbb{K}\subseteq\mathbb{Q}(\zeta_p)$. These modules are obtained from isomorphic subgroups corresponding to some special subspaces of $\mathbb{F}_p^s$, as presented below.

Consider ${\bf v}=(a_0,a_1,\ldots,a_{s-1})\in \mathbb{Z}^s$
and $\overline{\bf v}=(\overline{a_0},\overline{a_1},\ldots,\overline{a_{s-1}})$ in the vector space $V=\mathbb{F}_{p}^s$, where $\mathbb{F}_p$ is the finite field of $p$ elements, and where $\overline{a}\in\mathbb{F}_p$ denotes the reduction modulo $p$ of $a\in\mathbb{Z}$. Let \[ \mathcal{M}_{\bf v} = \{{\bf x}\in\mathbb{Z}^s:{\bf x}\cdot{\bf v}\equiv 0\pmod p\} = \mathcal{M}_{\bf v} = \{{\bf x}\in\mathbb{Z}^s:\overline{\bf x}\cdot\overline{\bf v}=\overline{0} \ \ \mbox{in} \ \ \mathbb{F}_p\}, \] where  $u\cdot v$ denotes the usual dot product between two vectors $u$ and $v$.

Since $\mathbb{F}_p^*$ is a cyclic group, there exists an element $\beta\in\mathbb{Z}$ such that $\beta$ is a $(p-1)$-th primitive root modulo $p$, that is, $\beta^{p-1}\equiv 1\pmod p$ and $\beta^j\not\equiv 1\pmod p$ for $j=1,2,\ldots,p-2$. The element $\alpha=\beta^{\frac{p-1}{s}}$ is a $s$-th primitive root modulo $p$. For each $i=0,1,\ldots,s-1$, denote \begin{equation} \label{eq-0} {\bf \mu_i}= \left(1,\alpha^i,\alpha^{2i},\ldots,\alpha^{(s-1)i}\right) \in \mathbb{Z}^s \end{equation} and \[ \overline{{\bf \mu_i}} = \left(1,\overline{\alpha}^i,\overline{\alpha}^{2i},\ldots,\overline{\alpha}^{(s-1)i}\right) \in \mathbb{F}_p^s. \] It is a known fact that the Vandermonde matrix \[ M := \left[ \begin{array}{cccc} 1 & 1 & \cdots & 1\\ 1 & \overline{\alpha} & \cdots & \overline{\alpha}^{(s-1)} \\ \vdots & \vdots & \ddots & \vdots \\ 1 & \overline{\alpha}^{s-1} & \cdots & (\overline{\alpha}^{s-1})^{s-1} \end{array} \right], \] has determinant
\[ \det(M) = \prod_{\substack{0\leq i<j<s}}(\overline{\alpha}^i-\overline{\alpha}^j), \] which is nonzero because
\begin{multline}\label{eq-1} \overline{\mu_i}\cdot\overline{\mu_j}  = \sum_{k=0}^{s-1} \overline{\alpha}^{k(i+j)} = \begin{cases} \frac{(\beta^{i+j})^s-1}{\beta^{i+j}-1} & \mbox{if } i+j\not\equiv 0\pmod s\\ s & \mbox{if } i+j\equiv 0\pmod s    \end{cases}\\ =  \begin{cases}
0 & \mbox{if } i+j\not\equiv 0\pmod s\\ s & \mbox{if } i+j\equiv 0\pmod s.\end{cases} \end{multline} Therefore, $\{\overline{\mu_0},\overline{\mu_1},\ldots,\overline{\mu_{s-1}}\}$ is a basis of $\mathbb{F}_p^s$.
If $s\equiv 1\pmod 2$, consider the subspace \[ W = \left\langle\overline{\mu_0},\overline{\mu_1},\ldots,\overline{\mu_{\frac{s-1}{2}}}\right\rangle_{\mathbb{F}_p} \subseteq\mathbb{F}_p^s. \]
Otherwise, if $s\equiv 0\pmod 2$, consider
\[  W=\left\langle\overline{\bf \mu_0},\overline{\bf \mu_1},\ldots,\overline{\bf \mu_{\frac{s}{2}}}\right\rangle_{\mathbb{F}_p} \subseteq\mathbb{F}_p^s. \] In each case, let \[ W^{\perp} = \{ {\bf x} \in \mathbb{F}_p^s:{\bf x}\cdot{\bf y}=0, \ \ \mbox{for all} \ \ {\bf y}\in W\} \] be the orthogonal complement of $W$.

\begin{proposition}\label{prop-2} If $s\equiv 1(mod\ 2)$, the orthogonal complement of $W$ is given by $$W^{\perp} = \left\langle\overline{\bf \mu_1},\overline{\bf \mu_2},\ldots,\overline{\bf \mu_{\frac{s-1}{2}}}\right\rangle_{\mathbb{F}_p} \subseteq\mathbb{F}_p^s.$$ \end{proposition}
\begin{proof} Let ${\bf x} = \displaystyle \sum_{i=0}^{s-1}a_i\overline{\bf \mu_i}\in \mathbb{F}_p^s$. Thus,  
\[ {\bf x}\in W^{\perp} \iff {\bf x}\cdot \overline{\bf \mu_j}=0, \ \ \mbox{for all} \ \ j=0,1,\ldots,\frac{s-1}{2}, \] that is, \[ \overline{x}\in W^{\perp} \iff \sum_{i=0}^{s-1} a_i\ \overline{\bf \mu_i}\cdot \overline{\bf\mu_j}=0, \ \ \mbox{for all} \ \ j=0,1,\ldots,\frac{s-1}{2}. \] Therefore, \[ {\bf x}\in W^{\perp}
\iff a_0s=a_{\frac{s+1}{2}}s = \ldots = a_{s-1}s =0, \] which implies that \[ a_i=0,~\mbox{for } i=0,a_{\frac{s+1}{2}},\ldots,a_{s-1}. \] So, ${\bf x}\in W^\perp$ if and only if ${\bf x}=a_1{\bf\mu_1}+\cdots+a_{\frac{s-1}{2}}{\bf \mu_{\frac{s-1}{2}}}$, which proves the result. \end{proof}

Similarly to Proposition \ref{prop-2}, we obtain the following result:

\begin{proposition}
    \label{cor-0} If $s\equiv 0\pmod 2$, the orthogonal complement of $W$ is given by $$W^{\perp} = \left<\overline{\bf \mu_1},\overline{\bf \mu_2},\ldots,\overline{\bf \mu_{\frac{s}{2}-1}}\right>.$$
\end{proposition} 

\subsection{The $\mathbb{Z}$-modules $\mathcal{M}_p$.}\label{subsection_Mp}

Since $\mathbb{K}$ is an algebraic number field of degree $s$, the $\mathbb{Z}$-modules $\mathcal{O}_\mathbb{K}$ and $\mathbb{Z}^s$ are isomorphic. More precisely, by the results of Section \ref{sec-2}, the map
\begin{equation}\label{eq_iso}
  \phi\left(a_0,a_1,\ldots,a_{s-1}\right)=\sum_{i=0}^{s-1} a_i\theta^i(t)  
\end{equation}
defines a $\mathbb{Z}$-module isomorphism from $\mathbb{Z}^s$ to $\mathcal{O}_\mathbb{K}$. Accordingly, for simplicity, we shall write ${\bf x}:=\phi({\bf x})\in\mathcal{O}_\mathbb{K}$ for all ${\bf x}=(a_0,a_1,\ldots,a_{s-1})\in\mathbb{Z}^s$.

Firstly, consider $s \equiv 1\pmod 2$. Let
\begin{equation}\label{M_p_s_1}
    \mathcal{M}_p:=\left\{{\bf x}\in\mathcal{O}_{\mathbb{K}}:{\bf x}\cdot {\bf \mu_i}\equiv 0\pmod p, \ \forall~i=0,1,\ldots,\frac{s-1}{2}\right\}
\end{equation}
be a $\mathbb{Z}$-submodule of $\mathcal{O}_{\mathbb{K}}$. As in the previous subsection, consider $$W = \left\langle\overline{\mu_0},\overline{\mu_1},\ldots,\overline{\mu_{\frac{s-1}{2}}}\right\rangle_{\mathbb{F}_p} \subseteq\mathbb{F}_p^s$$ and its orthogonal complement $W^{\perp} = \left\langle\overline{\bf \mu_1},\overline{\bf \mu_2},\ldots,\overline{\bf \mu_{\frac{s-1}{2}}}\right\rangle_{\mathbb{F}_p}$ (Proposition \ref{prop-2}).

If ${\bf x}\in\mathcal{M}_p$, then $\overline{\bf x}\in W^{\perp}$. Since $W^{\perp}\subseteq W$, it follows that $\overline{\bf x}\cdot\overline{\bf x}=\overline{0}$, that is, ${\bf x}\cdot {\bf x}\equiv 0\pmod p$. Furthermore, \[ {\bf x}\cdot \overline{\bf x} = \sum_{i=0}^{s-1}a_i^2 . \] If ${\bf x}\in W$, then ${\bf x}\cdot {\bf \mu_0}\equiv 0\pmod p$. Thus, \[ \sum_{i=0}^{s-1}a_i\equiv 0\pmod p, \ \ \mbox{implying} \ \ \left(\sum_{i=0}^{s-1}a_i\right)^2\equiv 0\pmod p^2. \] 
Since \[ Tr_{\mathbb{K}}({\bf x}\overline{\bf x}) = p\sum_{i=0}^{s-1}a_i^2 -\frac{p-1}{s}\left(\sum_{i=0}^{s-1}a_i\right)^2 ,\] it follows that \[ Tr_{\mathbb{K}}({\bf x}\overline{\bf x}) \equiv 0 \pmod {p^2}. \] If $t=\min\{Tr_{\mathbb{K}}({\bf x}\overline{\bf x}): {\bf x}\in\mathcal{M}_p, {\bf x}\neq 0\}$, then \[ \delta(\sigma(\mathcal{M}_p)) = \frac{t^{\frac{s}{2}}}{2^sp^\frac{s-1}{2}[\mathcal{O}_{\mathbb{K}}:\mathcal{M}_p]}. \] Now, if 
\[ \sum_{i=0}^{s-1}a_i\equiv 0\pmod 2, \ \ \mbox{then} \ \ \sum_{i=0}^{s-1}a_i^2\equiv 0 \pmod 2, \] and, therefore, \[ Tr_{\mathbb{K}}({\bf x}\overline{\bf x}) \equiv 0\pmod {2p^2}. \] Thus, \begin{equation} \label{eq-2} \delta(\sigma(\mathcal{M}_p)) \geq \frac{(2p^2)^{\frac{s}{2}}}{2^sp^{\frac{s-1}{2}}[\mathcal{O}_{\mathbb{K}}:\mathcal{M}_p]}. \end{equation}

Now, consider $s \equiv 0\pmod 2$. Let
\begin{equation} \label{M_p_s_0}
   \mathcal{M}_p := \{{\bf x}\in\mathcal{O}_{\mathbb{K}}:{\bf x}\cdot {\bf \mu}_i\equiv 0 \pmod{p}, \ \ \forall i=0,1,\ldots,s/2\} 
\end{equation}
\[  \] be a $\mathbb{Z}$-submodule of $\mathcal{O}_{\mathbb{K}}$.
As in the previous subsection, consider
\[  W=\left\langle\overline{\bf \mu_0},\overline{\bf \mu_1},\ldots,\overline{\bf \mu_{\frac{s}{2}}}\right\rangle_{\mathbb{F}_p} \subseteq\mathbb{F}_p^s \]
and its orthogonal complement $W^{\perp} = \{ {\bf x} \in \mathbb{F}_p^s:{\bf x}\cdot{\bf y}=0, \ \ \mbox{for all} \ \ {\bf y}\in W\}$.

If $\overline{\bf x}\in W^{\perp}$, then $\overline{\bf x}\cdot\overline{\bf x}=\overline{0}$, that is, $\displaystyle {\bf x}\cdot {\bf x}=\sum_{i=0}^{s-1}a_i^2\equiv 0 \pmod{p}$. Thus, \[ Tr_{\mathbb{K}}({\bf x}\cdot\overline{\bf x})\equiv 0 \pmod{p^2}. \] Therefore, \begin{equation}\label{eq-s-cong1}
    \delta(\sigma(\mathcal{M}_p)) \geq \frac{(p^2)^{\frac{s}{2}}}{2^sp^{\frac{s-1}{2}}[\mathcal{O}_{\mathbb{K}}:\mathcal{M}_p]}. \end{equation}

\subsection{The index $[\mathcal{O}_{\mathbb{K}}:\mathcal{M}_p]$}

Initially, in this subsection, we present some elementary results in algebra that are important for the subsequent computation of the index $[\mathcal{O}_{\mathbb{K}}:\mathcal{M}_p]$. Their proofs are omitted, since they follow from standard results in abstract algebra.

\begin{lemma} \label{lem-2} Let $N_1$ and $N_2$ be the submodules of a $\mathbb{Z}$-module $N$. Let $\pi:N_1\rightarrow \frac{N}{N_2}$ be a homomorphism given by $\pi(x)=x+N_2$, where $x\in N_1$. Then $\ker(\pi) = N_1\cap N_2$ and $N_1/(N_1\cap N_2)$ is isomorphic to $Im(\pi)\subseteq N/N_2$. Furthermore, if $[N:N_2]<\infty$, then $[N_1:N_1\cap N_2]<\infty$ is a divisor of $[N:N_2]$.\end{lemma}

\begin{lemma} \label{cor-2} For a prime number $p$, if $[N:N_1]=[N:N_2]=p$ and $N_1\neq N_2$, then $[N_1:N_1\cap N_2]=p$. \end{lemma} 

\begin{lemma} \label{lem-3} Let $N_1\subseteq N_2\subseteq N_3$ be $\mathbb{Z}$-modules. If $[N_3:N_2]$ and $[N_2:N_1]$ are finite, then $[N_3:N_1]=[N_3:N_2][N_2:N_1]$ is also finite. \end{lemma}

\begin{lemma} \label{lem-4} Let $a_1,a_2,\ldots,a_s\in\mathbb{Z}\setminus\{0\}$. If $d=\gcd(a_1,a_2,\ldots,a_s)$, then $d\mathbb{Z}=a_1\mathbb{Z}+\cdots+a_s\mathbb{Z}$. \end{lemma}

\begin{lemma} \label{cor-3} If $m\in\mathbb{Z}$ is a positive integer, then the map

\[\begin{array}{llll} \pi:&\mathbb{Z}^s&\rightarrow& \mathbb{Z}/m\mathbb{Z}\\
& (x_1,x_2,\ldots,x_s) &\mapsto & \overline{\sum_{i=1}^s x_i a_i}
\end{array} \]
is a homomorphism of $\mathbb{Z}$-modules. \end{lemma}

From Lemma \ref{cor-3} and the Isomorphism Theorem, it follows that \[ \mathbb{Z}^s/\ker(\pi) \simeq Im(\pi)\subseteq \mathbb{Z}/m\mathbb{Z}. \] Since $d=\gcd(a_1,a_2,\ldots,a_s)$, it follows from Lemma \ref{lem-4} that $\sum_{i=1}^s x_ia_i \in d\mathbb{Z}$. Thus, \[ Im(\pi) =\{\overline{0},\overline{d},\overline{2d},\ldots,\overline{(k-1)d}\}, \] where $kd$ is minimum positive integer divisible by $m$. Thus, $k=\frac{m}{gcd(m,d)}$. This implies that \[ [\mathbb{Z}^s:\ker(\pi)] = \frac{m}{gcd(m,d)}. \] If $\gcd(m,d)=1$, then $\pi$ is surjective and
\begin{multline}\label{eq-nova}
   \ker(\pi)=\left\{(x_1,x_2,\ldots,x_n)\in\mathbb{Z}^s:\displaystyle \sum_{i=1}^n x_ia_i\equiv 0 \pmod m\right\}\\
   = \{{\bf u}\in\mathbb{Z}^n:{\bf u}\cdot{\bf v}\equiv 0 \pmod m\},
\end{multline}
where ${\bf v}=(a_1,a_2,\ldots,a_n)\in\mathbb{Z}^n$. Hence, in this case, $\mathbb{Z}^n/\ker(\pi)$ is isomorphic to $\mathbb{Z}/m\mathbb{Z}$, which implies $[\mathbb{Z}^n:\ker(\pi)]=m$.
%Let  \[ \mathcal{M}_i=\{{\bf x}\in\mathbb{Z}^s:{\bf x}\cdot{\bf \mu_i}\equiv 0(mod\ p), \ \ \mbox{for} \ \ i=0,1,\ldots,s-1\}. \] for $i=0,1,\ldots,$
For each $i=0,1,\ldots,s-1$, consider ${\bf \mu_i}= \left(1,\alpha^i,\alpha^{2i},\ldots,\alpha^{(s-1)i}\right) \in \mathbb{Z}^s$ as in Equation \ref{eq-0}. Let \[ \mathcal{N}_i=\{{\bf x}\in\mathbb{Z}^s: {\bf x}\cdot{\bf \mu_i}\equiv 0 \pmod p\} \] be a $\mathbb{Z}$-submodule of $\mathbb{Z}^s$. Since $\gcd(1,\alpha^i,\alpha^{2i},\ldots,\alpha^{(s-1)i})$, it follows from (\ref{eq-nova}) that $$[\mathbb{Z}^s:\mathcal{N}_i]=p.$$

\begin{lemma} \label{lem-5} $\mathcal{N}_i\neq \mathcal{N}_j$ for each pair $i\neq j\in\{0,1,\ldots,s-1\}$. \end{lemma} \begin{proof} Suppose $i\neq j$, for $i,j=0,1,\ldots,s-1$. From (\ref{eq-1}), it follows that ${\bf \mu_{s-j}}\cdot{\bf \mu_j}=s$, thus ${\bf \mu_{s-j}}\notin \mathcal{N}_j$ for all $j$. Moreover, ${\bf \mu_{s-j}}\cdot{\bf \mu_i}=0$, which implies ${\bf \mu_{s-j}}\in\mathcal{N}_i$. Therefore, $\mathcal{N}_i\not\subseteq \mathcal{N}_j$. Analogously, $\mathcal{N}_j\not\subseteq \mathcal{N}_i$. Therefore, $\mathcal{N}_i\neq \mathcal{N}_j$. \end{proof}

For each $i=1,2,\ldots,s-2$, it follows from Lemma \ref{cor-2} that
\[[\mathcal{N}_0\cap\mathcal{N}_1\cap\cdots\cap\mathcal{N}_{i-1}:\mathcal{N}_0\cap\mathcal{N}_1\cap\cdots\cap\mathcal{N}_{i}]=p.  \]
From this and Lemma \ref{lem-3}, 
\begin{equation}\label{eq-nova2}
   [\mathbb{Z}^s:\mathcal{N}_0\cap\cdots\cap\mathcal{N}_{i}]=p^i,
\end{equation}
for $i=1,2,\ldots,s-1$.

Returning to the modules $\mathcal{M}_p$ defined in Subsection \ref{subsection_Mp}, the previous equation leads to the following theorem:


\begin{theorem} \label{prop-3} Consider the $\mathbb{Z}$-modules $\mathcal{M}_p$ defined in (\ref{M_p_s_1}) for the case $s\equiv 1\pmod 2$ and in (\ref{M_p_s_0}) for the case $s\equiv 0\pmod 2$. Thus, $$[\mathcal{O}_{\mathbb{K}}:\mathcal{M}_p]=\left\{\begin{matrix}
   p^{\frac{s+1}{2}}  &  \mbox{if}~s\equiv 1\pmod 2,\\
   p^{\frac{s+2}{2}}  &  \mbox{if}~s\equiv 0\pmod 2. \end{matrix}\right.$$ 
\end{theorem}

\begin{proof}
As pointed in (\ref{eq_iso}), $\mathbb{Z}^s$ is a $\mathbb{Z}$-module isomorphic to $\mathcal{O}_\mathbb{K}$. By (\ref{eq-nova}), $\mathcal{M}_p$ is isomorphic to $\ker(\pi)$ as a $\mathbb{Z}$-module. Hence, $\mathbb{Z}^s/\ker(\pi)\simeq \mathcal{O}_\mathbb{K}/\mathcal{M}_p$. Therefore, $[\mathcal{O}_{\mathbb{K}}:\mathcal{M}_p]=[\mathbb{Z}^s:\ker(\pi)]$. The desired result now follows from (\ref{eq-nova2}).
\end{proof}


\begin{thebibliography}{99}
\bibitem[samuel]{samuel} P. Samuel. \textit{Algebraic theory of numbers}, Hermann, Paris, 1982.
\bibitem[trajano]{trajano} T. P. N\'obrega Neto, J. C. Interlando and J. O. D. Lopes. \textit{On computing discriminats of subfields of $\mathbb{Q}(\zeta_{p^r})$}, J. Number Theory, 96, 319-325, 2002.
\bibitem[wash]{wash} L. C. Washington. \textit{Introduction to cyclotomic fields}, Springer New York, 1997.
\end{thebibliography}
\end{document}
